% =====================================================================
%  IEEE-STYLE CONFERENCE REPORT TEMPLATE  (self-contained)
%  MLE-AI-201 Machine Learning -- Final Project
%
%  This template MIMICS the IEEE two-column conference look using only
%  standard LaTeX packages. It does NOT depend on IEEEtran.cls, so it
%  compiles anywhere (tectonic, pdflatex, Overleaf) with nothing extra.
%
%  Build:   tectonic report.tex      (recommended)
%     or:   pdflatex report.tex ; pdflatex report.tex
% =====================================================================
\documentclass[10pt,twocolumn,a4paper]{article}

% ---- Geometry: IEEE-like margins / two narrow columns ----------------
\usepackage[a4paper,top=0.75in,bottom=1.0in,left=0.68in,right=0.68in,columnsep=0.25in]{geometry}

% ---- Times-like font (standard, ships with every TeX distro) ---------
\usepackage{mathptmx}          % Times text + math
\usepackage[T1]{fontenc}
\usepackage[utf8]{inputenc}    % allows accented / Azerbaijani chars

% ---- Common, dependency-free packages --------------------------------
\usepackage{graphicx}
\usepackage{amsmath,amssymb}
\usepackage{booktabs}
\usepackage{array}
\usepackage{titlesec}
\usepackage{caption}
\usepackage{cite}
\usepackage[hidelinks]{hyperref}
\usepackage{url}

% ---- IEEE-style section headings -------------------------------------
\renewcommand{\thesection}{\Roman{section}}
\renewcommand{\thesubsection}{\Alph{subsection}}
\titleformat{\section}[block]{\centering\normalsize\scshape}{\thesection.}{0.5em}{}
\titleformat{\subsection}[block]{\normalsize\itshape}{\thesubsection.}{0.5em}{}
\titlespacing*{\section}{0pt}{1.4ex plus 1ex minus .2ex}{1ex}
\titlespacing*{\subsection}{0pt}{1.0ex plus 1ex minus .2ex}{0.6ex}

% ---- Abstract / Index-Terms in the IEEE boldface-dash style ----------
\newcommand{\abstractkw}[1]{\noindent\textbf{\textit{#1}}}

\captionsetup{font=small,labelsep=period}
\setlength{\parindent}{1em}

% =====================================================================
\begin{document}

% ------------------------- TITLE BLOCK --------------------------------
\twocolumn[
  \begin{@twocolumnfalse}
  \begin{center}
    {\LARGE\bfseries Your Project Title Here:\\A Clear, Specific Subtitle\par}
    \vspace{1.2em}
    % Authors: Surname, Given name --- alphabetical by surname.
    {\large
      Aliyeva, Aysel\quad Mammadov, Murad\quad Rustamli, Rustam\par}
    \vspace{0.5em}
    {\itshape MLE-AI-201 Machine Learning, AI Academy --- Cohort II\par}
    \vspace{1.4em}
  \end{center}
  \begin{quote}
    \small
    \abstractkw{Abstract---}% IEEE uses an em-dash after "Abstract"
    One paragraph, roughly 150--250 words. State the problem (price
    prediction and the derived price-tier task on the bina.az dataset),
    what you built from scratch (a decision tree and an SVM), the key
    result (headline metrics vs.\ scikit-learn), and the main takeaway.
    Write it last. \textbf{End the abstract with your repository link},
    per IEEE convention:
    \href{https://github.com/your-team/your-repo}{\texttt{github.com/your-team/your-repo}}
    (tag \texttt{v1.0-final}).

    \vspace{0.6em}
    \abstractkw{Index Terms---}decision tree, support-vector machine,
    Pegasos, regression, classification, real-estate price prediction.
  \end{quote}
  \vspace{1.0em}
  \end{@twocolumnfalse}
]

% =====================================================================
\section{Introduction}
Motivate the problem and state your contributions as a short bulleted or
inline list. What is the task, why does it matter, and what did your
team do? Keep background brief; cite sources like this~\cite{pegasos}.

\section{Related Work / Background}
Briefly place your work in context: decision trees (impurity, information
gain), the soft-margin SVM and the hinge-loss primal, and the Pegasos
sub-gradient solver you used~\cite{pegasos}. One or two short paragraphs.

\section{Data and Preprocessing}
Describe the bina.az dataset: size, key columns, and the cleaning you
did (currency, types, missing values, duplicates, outliers). State the
two tasks (price regression; price-tier classification), how you derived
the tier label, and your train/validation/test split. \textbf{Name the
leakage columns you removed and why.}

\begin{table}[t]
  \centering
  \caption{Dataset summary after cleaning.}
  \label{tab:data}
  \begin{tabular}{@{}lr@{}}
    \toprule
    Property & Value \\
    \midrule
    Rows (after cleaning)       & --- \\
    Features used               & --- \\
    Train / Val / Test split    & --- / --- / --- \\
    Price-tier balance (\%)     & --- \\
    \bottomrule
  \end{tabular}
\end{table}

\section{Method}
\subsection{Decision tree (from scratch)}
Describe your impurity criteria (Gini / entropy / MSE reduction), the
split search, recursion, leaf rule, and hyperparameters.

\subsection{Support-vector machine (from scratch)}
State the objective and the Pegasos update you implemented:
\begin{equation}
  J(\mathbf{w},b)=\tfrac{\lambda}{2}\lVert\mathbf{w}\rVert^{2}
  +\tfrac{1}{n}\sum_{i=1}^{n}\max\!\big(0,\,1-y_i(\mathbf{w}^{\top}\mathbf{x}_i+b)\big).
\end{equation}
Explain the learning-rate schedule, feature scaling, and your non-linear
extension (kernelised Pegasos or random Fourier features), if any.

\section{Experimental Setup}
Metrics (RMSE/MAE/$R^2$ for regression; precision/recall/F1/AUC for
classification), hyperparameter search protocol, seeds, and the
scikit-learn baselines you compared against (with matched settings).

\section{Results}
Lead with a results table; back every claim with a number or figure.

\begin{table}[t]
  \centering
  \caption{Headline results (test set, mean $\pm$ std over folds).}
  \label{tab:results}
  \begin{tabular}{@{}lcc@{}}
    \toprule
    Model & Task A (RMSE) & Task B (F1) \\
    \midrule
    Ours: decision tree     & --- & --- \\
    Ours: SVM (Pegasos)     & n/a & --- \\
    scikit-learn baseline   & --- & --- \\
    \bottomrule
  \end{tabular}
\end{table}

\begin{figure}[t]
  \centering
  % \includegraphics[width=\linewidth]{figures/your_plot.pdf}
  \fbox{\parbox[c][3.2cm][c]{0.9\linewidth}{\centering\small
    Replace with a figure (e.g.\ validation curve, confusion matrix,
    or feature importances). Save plots to \texttt{figures/}.}}
  \caption{One clear figure beats ten decorative ones.}
  \label{fig:main}
\end{figure}

\section{Discussion and Limitations}
Where and why each model fails (luxury outliers, sparse districts,
scaling effects), the bias--variance story, and honest limitations.

\section{Conclusion}
One short paragraph: what you found and what you would do next.

\section*{Tools and Acknowledgements}
Disclose any AI-assistant use and what it was used for, plus data/code
you built on. (Required.)

% =====================================================================
\begin{thebibliography}{9}
\bibitem{pegasos}
S.~Shalev-Shwartz, Y.~Singer, N.~Srebro, and A.~Cotter,
``Pegasos: primal estimated sub-gradient solver for SVM,''
\emph{Mathematical Programming}, vol.~127, no.~1, pp.~3--30, 2011.
\end{thebibliography}

\end{document}
